%%
%% Author: Deeps
%% 22.06.2018
%%

% Preamble
\documentclass[
a4paper,		% Verwenden von A4 als Seitenformat
12pt,			% Festlegung der Standardtextgröße auf 12pt
pagesize,
headsepline,
titlepage		% Erstellen einer Titelseite, Automatisches zentrieren der Titelseite
]{scrartcl}

% Packages
\usepackage[ngerman]{babel}		% deutsche Trennmuster
\usepackage[utf8]{inputenc}		% direkte Eingabe von Umlauten & Co. (Vorsicht: Encoding im Editor muss auch UTF-8 sein!)

\usepackage[T1]{fontenc}			% T1-Schriften

\usepackage{mathptmx}			% Times/Mathe \rmdefault
\usepackage[scaled=.90]{helvet}	% Skalierte Helvetica \sfdefault
\usepackage{courier}			% Courier \ttdefault

\usepackage{
amsmath,	% Mehr mathematische Symbole
amsthm,		% Mehr mathematische Symbole
amsfonts,	% Mehr mathematische Symbole
graphicx, 	% Einbindung von Grafiken
caption 		% Bildunterschriften
}

% Wenn man direkt mit dem pdflatex eine PDF-Datei erzeugt, sollten diese beiden Pakete eingebunden werden
\usepackage{hyperref} % Hyperlinks anklickbar
\usepackage{ae, aecompl} % bessere Bildschirmschriftarten

% Code anzeigen
\usepackage{listings}
\usepackage{algorithm}
\usepackage{algpseudocode}

\usepackage{natbib}

\usepackage{url}        %bibtex
\usepackage{graphicx}

% Styling
\pagestyle{headings}

\headsep4mm % Abstand der Kopfzeile vom Text:

\typearea[current]{current}     % Satzspiegel neu berechnen

% Title

\titlehead{
\hfill \includegraphics[height=3cm]{resources/up_logo.eps}
\hfill \includegraphics[height=3cm]{resources/hpi_logo.eps}
}
\subject{Bachelorarbeit}
\title{
Maschinelles Lernen im Onlinehandel: \\
Extraktion produktspezifischer Daten
\vspace{1.25cm}
}
\author{
Leonardo Hübscher\\
\bigskip
\\B.Sc.\\
IT-Systems Engineering\\
\bigskip
Fachgebiet für Informationssysteme
\vspace{0.75cm}
\\Betreuer:\\
Prof. Felix Naumann\\
Leon Bornemann\\
}
\date{20. Juli 2018}

% Document
\begin{document}

    \maketitle

    \begin{abstract}
        \section*{Zusammenfassung}
\markboth{Zusammenfassung}{Zusammenfassung}
\label{sec:abstract}

Durch die Vielzahl von Onlineshops und Fülle an Angeboten verliert der Onlinekäufer schnell die Übersicht.
Preisvergleichsplattformen wie idealo helfen dem Kunden das günstigste Angebot im Netz zu finden.
Die Gewährleistung der möglichst vollständigen Markttransparenz ist eine grundlegende Herausforderung für idealo.
Das von uns entwickelte Softwaresystem \textit{Scout} soll dabei helfen, den Produktkatalog von idealo auf
Vollständigkeit zu überprüfen und fehlende Angebote aufzulisten.
Ein wichtiger Prozessschritt ist dabei die Extrahierung von Produktinformationen, wie Produktname oder Preis, aus den
einzelnen Webseiten.
Die Schwierigkeit der Extraktion liegt darin, dass jeder Shop einen individuellen Aufbau besitzt und unterschiedlich
strukturiert ist.

Das entwickelte Parser-Modul löst dieses Problem, indem es für jeden Shop eigene Regeln zur Extraktion der
Produktinformationen verwendet.
Dabei ist es nicht erforderlich, dass diese Regeln manuell erstellt werden müssen.
Durch die Nutzung der bereits vorhandenen Angebote aus dem Bestand von idealo kann die Extrahierung der Struktur
mittels maschinellen Lernens erfasst werden.
Messungen, welche auf 50 verschiedenen Shops basieren, haben ergeben, dass die Produktinformationen mit einer
Precision von über 95 Prozent bei einer Accuracy von etwa 50\% extrahiert werden können.
    \end{abstract}

    \tableofcontents
    \newpage

    \bibliography{main}
    \bibliographystyle{plaindin}      % BibTeX Styles nach Norm DIN 1505

\end{document}
